% ==============================================================
%  Exercices de révision MTH8304 — Corrigé complet
%  Source : fichier photos.pdf (6 pages, notes manuscrites)
%  Sujets : classifieur probabiliste (Bishop 6.4.9),
%           classification d'écoulement (Partie II),
%           modélisation de sédiment (distributions)
%  Compilable sur Overleaf avec pdfLaTeX
% ==============================================================
\documentclass[12pt,a4paper]{article}

\usepackage[utf8]{inputenc}
\usepackage[T1]{fontenc}
\usepackage[french]{babel}
\usepackage[top=2.5cm,bottom=2.5cm,left=2.8cm,right=2.8cm]{geometry}
\usepackage{amsmath,amssymb,amsthm}
\usepackage{bm}
\usepackage{xcolor}
\usepackage{tcolorbox}
\usepackage{booktabs}
\usepackage{enumitem}
\usepackage{lmodern}
\usepackage{microtype}
\usepackage{titlesec}
\usepackage{mdframed}

\tcbuselibrary{skins,breakable,theorems}

% ---- couleurs ------------------------------------------------
\definecolor{bleu}{RGB}{20,90,170}
\definecolor{vert}{RGB}{0,120,60}
\definecolor{rouge}{RGB}{180,30,30}
\definecolor{gris}{RGB}{245,245,248}
\definecolor{orange}{RGB}{200,100,0}

% ---- boîtes --------------------------------------------------
\newtcolorbox{enonce}[1]{
  colback=bleu!8, colframe=bleu, fonttitle=\bfseries,
  title=Énoncé — #1, breakable}

\newtcolorbox{solution}{
  colback=vert!5, colframe=vert, fonttitle=\bfseries,
  title=Solution détaillée, breakable}

\newtcolorbox{rappel}[1]{
  colback=gris, colframe=gray!50, fonttitle=\bfseries\small,
  title=Rappel : #1, breakable}

% ---- macros --------------------------------------------------
\newcommand{\T}{^{\!\top}}
\newcommand{\R}{\mathbb{R}}
\newcommand{\E}{\mathbb{E}}
\newcommand{\Var}{\operatorname{Var}}

% ---- sections ------------------------------------------------
\titleformat{\section}{\large\bfseries\color{bleu}}{\thesection.}{0.6em}{}[\vspace{-4pt}\rule{\linewidth}{0.8pt}]
\titleformat{\subsection}{\normalsize\bfseries\color{bleu}}{\thesubsection}{0.5em}{}

\setlength{\parskip}{6pt}
\setlength{\parindent}{0pt}

% ==============================================================
\begin{document}

\begin{center}
  {\LARGE\bfseries\color{bleu} Exercices de révision — MTH8304}\\[6pt]
  {\large Classifieur probabiliste \textbullet{} Classification \textbullet{} Modélisation}\\[4pt]
  {\small Polytechnique Montréal \quad|\quad Corrigé complet}
\end{center}

\vspace{8pt}

% ==============================================================
\section{Exercice 1 — Bishop 6.4.9 : Classifieur probabiliste pour $K$ classes}
% ==============================================================

\begin{enonce}{Contexte et notations}
On considère un classifieur probabiliste avec $K$ classes $\mathcal{C}_1, \ldots, \mathcal{C}_K$ :
\begin{itemize}[itemsep=2pt]
  \item $P(\mathcal{C}_k) = \Pi_k$ : probabilité \textbf{a priori} de la classe $k$
  \item $P(\phi \mid \mathcal{C}_k)$ : probabilité conditionnelle de classe ($\phi$ = vecteur de caractéristiques)
  \item $\{(\phi_n, t_n)\}_{n=1}^{N}$ : observations i.i.d., avec $t_n \in \{0,1\}^K$ en encodage \emph{one-hot}
  \item $t_{nk} = 1$ si l'observation $n$ appartient à la classe $k$, $0$ sinon
\end{itemize}
\textbf{Question :} Montrez que l'estimateur par maximum de vraisemblance de $\Pi_k$ est
\[
  \hat{\Pi}_k = \frac{N_k}{N}
\]
où $N_k = \sum_{n=1}^{N} t_{nk}$ est le nombre d'observations de la classe $k$.
\end{enonce}

\begin{solution}

\textbf{Étape 1 : Probabilité d'une observation $(\phi_n, t_n)$.}

La loi jointe s'écrit :
\[
  P(\phi_n, t_n)
  = \prod_{k=1}^{K} P(\phi_n, \mathcal{C}_k)^{t_{nk}}
  = \prod_{k=1}^{K} \bigl[P(\phi_n \mid \mathcal{C}_k)\,\Pi_k\bigr]^{t_{nk}}.
\]
(Seul le facteur $k$ tel que $t_{nk}=1$ contribue réellement au produit.)

\medskip
\textbf{Étape 2 : Vraisemblance totale (loi multinomiale, $N$ observations indépendantes).}
\[
  \mathcal{L}(\Pi_1,\ldots,\Pi_K)
  = \prod_{n=1}^{N} P(\phi_n, t_n)
  = \prod_{n=1}^{N}\prod_{k=1}^{K} \bigl[P(\phi_n \mid \mathcal{C}_k)\,\Pi_k\bigr]^{t_{nk}}.
\]

\medskip
\textbf{Étape 3 : Log-vraisemblance.}
\[
  \ln \mathcal{L}
  = \sum_{n=1}^{N}\sum_{k=1}^{K} t_{nk}
    \bigl[\ln P(\phi_n \mid \mathcal{C}_k) + \ln \Pi_k\bigr].
\]
On sépare les termes dépendant de $\Pi_k$ :
\[
  \ln \mathcal{L}
  = \underbrace{\sum_{n,k} t_{nk}\ln P(\phi_n \mid \mathcal{C}_k)}_{\text{indépendant de }\Pi_k}
  + \sum_{k=1}^{K} N_k \ln \Pi_k,
\]
avec $N_k = \sum_{n=1}^{N} t_{nk}$.

\medskip
\textbf{Étape 4 : Maximisation sous contrainte $\sum_{k=1}^{K}\Pi_k = 1$ (multiplicateur de Lagrange).}

On forme le Lagrangien :
\[
  \mathcal{F}(\Pi_1,\ldots,\Pi_K,\lambda)
  = \sum_{k=1}^{K} N_k \ln \Pi_k
    + \lambda\!\left(\sum_{k=1}^{K} \Pi_k - 1\right).
\]

\medskip
\textbf{Étape 5 : Dérivée par rapport à $\Pi_k$.}
\[
  \frac{\partial \mathcal{F}}{\partial \Pi_k}
  = \frac{N_k}{\Pi_k} + \lambda = 0
  \quad\Longrightarrow\quad
  \Pi_k = -\frac{N_k}{\lambda}.
\]

\medskip
\textbf{Étape 6 : Détermination de $\lambda$ via la contrainte.}
\[
  \sum_{k=1}^{K} \Pi_k = 1
  \quad\Rightarrow\quad
  \sum_{k=1}^{K} \left(-\frac{N_k}{\lambda}\right) = 1
  \quad\Rightarrow\quad
  -\frac{N}{\lambda} = 1
  \quad\Rightarrow\quad
  \lambda = -N.
\]

\medskip
\textbf{Conclusion :}
\[
  \boxed{\hat{\Pi}_k = \frac{N_k}{N}}
\]
L'estimateur MLE de la probabilité a priori est simplement la fréquence observée de la classe $k$.
\end{solution}

% ==============================================================
\section{Partie II — Classification des mesures d'écoulement d'une rivière}
% ==============================================================

\begin{enonce}{Contexte}
On classe les mesures d'écoulement d'une rivière en deux catégories :
\textbf{sous-critique} ($\mathrm{Fr} < 1$) et \textbf{sur-critique} ($\mathrm{Fr} > 1$).
\begin{itemize}[itemsep=2pt]
  \item Données d'entrée : vitesse $v$ et profondeur $d$
  \item 300 observations : \textbf{270 sous-critique} (90\,\%) et \textbf{30 sur-critique} (10\,\%)
  \item Performance du modèle : \textbf{95\,\% en entraînement}, \textbf{61\,\% en test}
\end{itemize}
\end{enonce}

% -----------------------------------------------------------
\subsection{Question (a) — Deux explications de la différence train/test}
% -----------------------------------------------------------

\begin{enonce}{Q(a)}
Donnez 2 explications différentes pour expliquer la différence de performance entre l'entraînement (95\,\%) et le test (61\,\%).
\begin{enumerate}[label=\roman*)]
  \item Décrivez le mécanisme causant cette différence.
  \item Commentez sur l'interprétation des mesures de performance.
  \item Proposez un diagnostic pour vérifier quelle explication est en cause.
\end{enumerate}
\end{enonce}

\begin{solution}

\textbf{i) Deux mécanismes possibles :}

\begin{enumerate}[label=\textcircled{\small\arabic*},itemsep=4pt]
  \item \textbf{Déséquilibre de classes.}
  Avec 90\,\% de données sous-critiques, un modèle ``trivial'' qui prédit toujours
  \emph{sous-critique} obtient déjà 90\,\% d'accuracy en entraînement.
  Si la proportion change en test (par exemple plus d'événements sur-critiques),
  l'accuracy chute brutalement.

  \item \textbf{Sur-apprentissage (\emph{overfitting}).}
  Le modèle mémorise les données d'entraînement (bruit inclus) sans apprendre
  la structure générale. Il performe bien sur le train mais généralise mal sur
  des données non vues.
\end{enumerate}

\medskip
\textbf{ii) Interprétation des mesures de performance :}

L'accuracy est une mesure \textbf{trompeuse} sur des données déséquilibrées.
Avec 90\,\% de sous-critiques, prédire systématiquement la classe majoritaire
donne 90\,\% d'accuracy sans apprendre quoi que ce soit.
Il faudrait utiliser des métriques plus robustes : précision, rappel, score F1,
matrice de confusion, ou AUC-ROC.

\medskip
\textbf{iii) Diagnostics :}

\begin{itemize}[itemsep=4pt]
  \item Pour le \textbf{déséquilibre} : vérifier les proportions de chaque classe
  \emph{prédite} en test. Si le modèle ne prédit jamais ``sur-critique'', c'est
  le déséquilibre qui est en cause.

  \item Pour le \textbf{sur-apprentissage} : tracer les courbes d'erreur
  entraînement \emph{et} test en fonction de la taille de l'ensemble
  d'entraînement (\emph{learning curves}).
  Un grand écart entre les deux courbes révèle un sur-apprentissage.
\end{itemize}
\end{solution}

% -----------------------------------------------------------
\subsection{Question (b) — Nombre de Froude et régression logistique}
% -----------------------------------------------------------

\begin{enonce}{Q(b)}
Le nombre de Froude est défini par
\[
  \mathrm{Fr} = \frac{v}{\sqrt{g\,d}}\,,
\]
où $g$ est l'accélération gravitationnelle.
On utilise la \textbf{régression logistique} comme classifieur dans l'espace $(v, d)$.
Expliquez en quoi ce modèle est \textbf{mal adapté} pour cette classification.
Proposez des solutions.
\end{enonce}

\begin{solution}

\textbf{Problème fondamental : frontière de décision non-linéaire.}

La régression logistique est un \textbf{classifieur linéaire} : sa frontière de décision
dans l'espace $(v, d)$ est un hyperplan de la forme
\[
  w_1 v + w_2 d + b = 0.
\]

Or la frontière naturelle entre sous-critique et sur-critique est $\mathrm{Fr} = 1$, soit :
\[
  \frac{v}{\sqrt{g\,d}} = 1
  \quad\Longleftrightarrow\quad
  v = \sqrt{g\,d}
  \quad\Longleftrightarrow\quad
  v^2 = g\,d.
\]
C'est une \textbf{courbe parabolique} dans l'espace $(v, d)$ : non-linéaire.
La régression logistique ne peut donc pas capturer cette frontière.

\medskip
\textbf{Solutions :}

\begin{enumerate}[label=\textbf{--},itemsep=4pt]
  \item \textbf{Classifier directement par rapport à $\mathrm{Fr}$} (une seule feature) :
  la frontière $\mathrm{Fr} = 1$ devient triviale à apprendre avec n'importe quel classifieur.

  \item \textbf{Classifier dans l'espace $(v,\,\sqrt{d})$} (ou $(v,\,\sqrt{g\,d})$) :
  la frontière $v = \sqrt{g\,d}$ devient linéaire dans ce nouvel espace,
  et la régression logistique devient appropriée.
\end{enumerate}
\end{solution}

% ==============================================================
\section{Exercice 2 — Modélisation de sédiment en suspension}
% ==============================================================

\begin{enonce}{Contexte}
On dispose de $N = 200$ mesures journalières de sédiment en suspension dans une rivière.
\begin{itemize}[itemsep=2pt]
  \item Les données semblent \textbf{unimodales}.
  \item La distribution est \textbf{asymétrique vers la droite} (\emph{right-skewed}).
\end{itemize}
\end{enonce}

% -----------------------------------------------------------
\subsection{Question (a) — Gaussienne ajustée par maximum de vraisemblance}
% -----------------------------------------------------------

\begin{enonce}{Q(a-b) — Impact des événements extrêmes}
On ajuste une gaussienne $\mathcal{N}(\mu, \sigma^2)$ par maximum de vraisemblance.

Expliquez pourquoi et comment la présence d'un \textbf{petit nombre d'événements extrêmes}
affecte l'estimation de la variance.
\end{enonce}

\begin{solution}

L'estimateur MLE de la variance est :
\[
  \hat{\sigma}^2_{\mathrm{MLE}}
  = \frac{1}{N}\sum_{n=1}^{N}(x_n - \hat{\mu})^2.
\]

Les événements extrêmes (grandes valeurs $x_n \gg \hat{\mu}$) contribuent
\textbf{quadratiquement} au terme $(x_n - \hat{\mu})^2$.
Même un petit nombre d'extrêmes peut donc \textbf{dominer} la somme et
\textbf{surestimer artificiellement} $\hat{\sigma}^2$.

La gaussienne ajustée sera trop étalée (\emph{over-dispersed}),
ce qui correspond mal à la réalité des données typiques.
\end{solution}

\begin{enonce}{Q(a-c) — Conséquences pratiques pour les valeurs typiques}
Quelles sont les conséquences pratiques si on cherche à estimer les probabilités
des valeurs \emph{typiques} (autour de $\mu$) ?
\end{enonce}

\begin{solution}

La densité gaussienne est normalisée : $\int_{-\infty}^{+\infty} f(x)\,dx = 1$.
Si $\hat{\sigma}^2$ est surestimé, la hauteur du pic (densité maximale en $\mu$) :
\[
  f(\mu) = \frac{1}{\sqrt{2\pi\,\hat{\sigma}^2}}
\]
est \textbf{réduite}. Les probabilités des valeurs typiques (proches de $\mu$)
sont ainsi \textbf{sous-estimées}, au profit des queues.

\textbf{Conséquence pratique :}
Si l'on cherche à prédire la probabilité d'une mesure sédiment normale (événement courant),
le modèle gaussien calibré sur des extrêmes la sous-estimera systématiquement,
conduisant à une mauvaise gestion du risque.
\end{solution}

% -----------------------------------------------------------
\subsection{Question (b) — Choix parmi trois modèles}
% -----------------------------------------------------------

\begin{enonce}{Q(b)}
On dispose de trois modèles alternatifs :
\begin{enumerate}[label=\arabic*.]
  \item Gaussienne $\mathcal{N}(\mu, \sigma^2)$
  \item Loi de Student $t_\nu(\mu, \sigma^2)$
  \item Mélange de $M=5$ gaussiennes
\end{enumerate}
Expliquez lequel des trois modèles vous choisissez et pourquoi, en vous appuyant sur :
\begin{enumerate}[label=(\roman*)]
  \item L'adéquation du modèle à ce que vous connaissez des données.
  \item Le risque de sur-ajustement étant donné la taille d'entraînement ($N=200$).
  \item La facilité à bien ajuster chaque modèle.
\end{enumerate}
\end{enonce}

\begin{solution}

\begin{center}
\renewcommand{\arraystretch}{1.8}
\begin{tabular}{@{}lccc@{}}
\toprule
\textbf{Critère} & \textbf{Gaussienne} & \textbf{$t$ de Student} & \textbf{Mélange 5G} \\
\midrule
(i) Adéquation aux données & Faible & \textbf{Bonne} & Excessive \\
(ii) Risque sur-ajustement & Faible (2 params) & \textbf{Faible (3 params)} & Élevé (14 params) \\
(iii) Facilité d'ajustement & Analytique & Simple & Complexe (EM) \\
\bottomrule
\end{tabular}
\end{center}

\medskip
\textbf{Modèle choisi : loi de Student $t_\nu(\mu, \sigma^2)$}

\begin{enumerate}[label=(\roman*),itemsep=6pt]
  \item \textbf{Adéquation :}
  Le $t$ de Student possède des \textbf{queues lourdes} (heavy-tailed), contrôlées par
  le paramètre de degrés de liberté $\nu$.
  Cela permet de mieux absorber les événements extrêmes sans qu'ils contaminent
  l'estimation de la dispersion centrale.
  La gaussienne ($\nu \to \infty$) est trop légère en queue ; le mélange de 5 gaussiennes
  permet une multimodalité inutile pour des données \emph{unimodales}.

  \item \textbf{Risque de sur-ajustement :}
  Le $t$ de Student n'a que \textbf{3 paramètres} $(\mu, \sigma, \nu)$
  pour $N = 200$ observations → risque très faible.
  Le mélange de 5 gaussiennes nécessite jusqu'à $5 \times 3 - 1 = 14$ paramètres libres,
  ce qui est disproportionné pour 200 données.

  \item \textbf{Facilité d'ajustement :}
  L'estimation des paramètres du $t$ de Student se fait par une optimisation numérique
  simple (algorithme EM ou MLE direct), beaucoup plus stable que l'EM
  du mélange de gaussiennes (initialisation sensible, convergence non garantie).
\end{enumerate}

\textbf{Conclusion :} Le $t$ de Student offre le meilleur compromis entre flexibilité
(queues lourdes adaptées aux extrêmes), parcimonie (3 paramètres) et facilité d'ajustement.
\end{solution}

% ==============================================================
%  RÉCAPITULATIF
% ==============================================================
\newpage
\section*{Récapitulatif des résultats clés}

\begin{tcolorbox}[colback=bleu!8,colframe=bleu,fonttitle=\bfseries,
  title={Exercice 1 — Bishop 6.4.9}]
Estimateur MLE des probabilités a priori d'un classifieur probabiliste à $K$ classes :
\[
  \hat{\Pi}_k = \frac{N_k}{N}
\]
Dérivé par maximisation du Lagrangien $\sum_k N_k \ln \Pi_k + \lambda\!\left(\sum_k \Pi_k - 1\right)$.
\end{tcolorbox}

\vspace{8pt}

\begin{tcolorbox}[colback=bleu!8,colframe=bleu,fonttitle=\bfseries,
  title={Partie II — Classification d'écoulement}]
\begin{itemize}[itemsep=4pt]
  \item Performance train/test : deux causes possibles =
        \textbf{déséquilibre de classes} ou \textbf{sur-apprentissage}
  \item Diagnostic : \emph{learning curves} + proportions de classes prédites
  \item Régression logistique inadaptée car $\mathrm{Fr}=v/\!\sqrt{gd}=1$
        est une frontière \textbf{non-linéaire} dans $(v,d)$
  \item Solutions : classer sur $\mathrm{Fr}$ directement, ou sur $(v, \sqrt{d})$
\end{itemize}
\end{tcolorbox}

\vspace{8pt}

\begin{tcolorbox}[colback=bleu!8,colframe=bleu,fonttitle=\bfseries,
  title={Exercice 2 — Modélisation de sédiment}]
\begin{itemize}[itemsep=4pt]
  \item Les événements extrêmes \textbf{surestiment} $\hat{\sigma}^2_{\mathrm{MLE}}$ de la gaussienne
        (contribution quadratique), entraînant une \textbf{sous-estimation} des probabilités typiques
  \item Modèle choisi : \textbf{$t$ de Student} — queues lourdes + 3 params + ajustement simple
\end{itemize}
\end{tcolorbox}

\end{document}
